% titel: Blatt 2, A4: Verstaendnisfrage -- Arbeit und Leistung
% skript: Kap. Energie, Arbeit, Leistung

% Kinetische Energie, Momentanleistung.

Arbeit und Leistung verbinden Kraft, Weg und Zeit; die kinetische Energie
ist an die Masse und das Quadrat der Geschwindigkeit gekoppelt.

\begin{werkzeugkasten}[Werkzeuge]
\textbf{Kinetische Energie}
\begin{equation}
  E_{\mathrm{kin}}=\tfrac12 m v^2.
\end{equation}

\textbf{Momentanleistung} (Kraft $\times$ Geschwindigkeit)
\begin{equation}
  P=\frac{\mathrm dW}{\mathrm dt}=F\,v=m\,a\,v\;.
\end{equation}
Bei gleichmaessiger Beschleunigung aus der Ruhe gilt
$v=at$, $s=\tfrac12 at^2$.

\textbf{Arbeit.} Die von einer Kraft verrichtete Arbeit ist die
Energieubertragung: $W=\Delta E_{\mathrm{kin}}$.
\end{werkzeugkasten}

\begin{werkzeugkasten}[Vorgehen]
\begin{itemize}
  \item Gleiche kinetische Energie: $\tfrac12 m v^2$ gleichsetzen und nach
        $v$ aufloesen.
  \item Momentanleistung zu einem Zeitpunkt: $P(t)=m\,a\,v(t)$, mit
        $v(t)=at$.
  \item Die Beschleunigung aus $v(t)$ und gemessener Zeit gewinnen
        ($v=v_0+at$, hier $v_0=0$).
\end{itemize}
\end{werkzeugkasten}

\begin{achtung}[Typische Fallen]
\begin{itemize}
  \item Die Geschwindigkeit geht quadratisch ein: ``doppelte Masse bei
        gleicher Energie'' heisst nicht ``doppelte Geschwindigkeit''.
  \item Momentanleistung braucht die \emph{momentane} Geschwindigkeit,
        nicht die Endgeschwindigkeit.
  \item Masse $m$ nicht vergessen ($P=mav$, nicht $P=av$).
\end{itemize}
\end{achtung}
